% sections/09_h1_weather.tex — V5a-V5d + H1 verdict (v3 post-outlier numbers)

\section{H1: Does Adverse Weather Increase Accident Severity?}

Three complementary visualisations address this hypothesis. All figures reflect
the cleaned dataset (2016--2022, $2\times$IQR outlier removal).

%% ── V5a ─────────────────────────────────────────────────────
\subsection{V5a: Weather Group vs \% of Accident Counts --- 100\% Stacked Bar}

Each weather group is shown as a 100\% stacked bar, split between High (3--4)
and Low (1--2) severity. This normalises for volume differences between weather
categories, enabling direct proportion comparison.

% \figchart{v5a_weather_stacked}{Weather group vs high-severity proportion}{fig:v5a}
\begin{figure}[H]
  \centering
  \figchart{v5a_weather_stacked}{Weather group vs.\% of Accident Counts (100\% stacked bar)}{fig:v5a}
  \label{fig:v5a}
\end{figure}

\idiomtable
  {1D Bar (normalised stack)}
  {Vertical Length (proportion) + Colour Hue (High = dark, Low = light)}
  {Compare proportions across weather categories}
  {7 weather groups --- optimal for grouped comparison}

\insightbox{After outlier removal, Rain carries the highest high-severity
proportion (24.38\%), followed by Cloudy (22.43\%) and Snow/Ice (21.45\%).
Notably, Snow/Ice has moved much closer to Rain after cleaning, narrowing the
gap from 3.3\,pp to just 2.9\,pp. Clear weather (17.85\%) still ranks above
Fog/Haze (16.93\%) and Wind/Dust (13.77\%), reinforcing the severity-definition
effect: fog accidents clear quickly despite perceived danger. The cleaner
dataset strengthens the Rain--severity signal.}

%% ── V5b ─────────────────────────────────────────────────────
\subsection{V5b: Weather Group vs \% of Total Accidents --- Absolute Share}

This stacked bar chart shows each weather group's absolute share of total
accidents, split by severity.

% \figchart{v5b_weather_absolute}{Absolute share of accidents by weather group}{fig:v5b}
\begin{figure}[H]
  \centering
  \figchart{v5b_weather_absolute}{Absolute share of total accidents by weather group and severity}{fig:v5b}
  \label{fig:v5b}
\end{figure}

\insightbox{While rain and adverse weather show higher severity \emph{rates},
they represent a small fraction of total accidents (Rain: $\approx$8\%,
Snow/Ice: $\approx$4\%). The overwhelming majority of high-severity accidents
in absolute terms occur in Clear and Cloudy conditions, simply because these
dominate overall volume.}

\begin{quote}
  \itshape
  This duality --- higher rate in bad weather, higher absolute count in good
  weather --- is crucial for policy interpretation.
  \textbf{Charts V5a and V5b must always be read together, never independently.}
\end{quote}

%% ── V5c ─────────────────────────────────────────────────────
\subsection{V5c: Weather Group vs High Severity Ratio \% --- Ranked Bar}

A sorted horizontal bar chart ranks all 7 weather groups by their high-severity
ratio with distinct colours per category.

% \figchart{v5c_weather_ranked}{Weather groups ranked by high-severity ratio}{fig:v5c}
\begin{figure}[H]
  \centering
 \figchart{v5c_weather_ranked}{Weather groups ranked by high-severity ratio}{fig:v5c}
  \label{fig:v5c}
\end{figure}

\idiomtable
  {1D Bar}
  {Horizontal Length (ordered, aligned) + Colour Hue (categorical identity)}
  {Rank and compare severity rate across weather types}
  {7 categories --- direct visual ranking}

\begin{table}[H]
\centering\small
\begin{tabular}{@{}lcc@{}}
  \toprule
  \hdrrow \textbf{Weather Group} & \textbf{High Sev.\ Rate} & \textbf{Rank} \\ \midrule
  \negrow Rain      & 24.38\% & 1 \\
  \negrow Cloudy    & 22.43\% & 2 \\
  \midrow Snow/Ice  & 21.45\% & 3 \\
  \midrow Clear     & 17.85\% & 4 \\
  \posrow Fog/Haze  & 16.93\% & 5 \\
  \posrow Wind/Dust & 13.77\% & 6 \\
  \posrow Other     & 11.38\% & 7 \\
  \bottomrule
\end{tabular}
\caption{Weather group high-severity rates (post-$2\times$IQR removal, 2016--2022)}
\end{table}

%% ── V5d ─────────────────────────────────────────────────────
\subsection{V5d: Visibility Bins vs High Severity \%}

Accidents are binned by visibility into five ranges and compared by severity
proportion.

% \figchart{v5d_visibility}{Visibility bins vs high-severity proportion}{fig:v5d}
\begin{figure}[H]
  \centering
  \figchart{v5d_visibility}{Visibility bins vs high-severity proportion}{fig:v5d}
  \caption{Visibility bins vs.\ high-severity proportion}
  \label{fig:v5d}
\end{figure}

\idiomtable
  {1D Bar (normalised stack)}
  {Vertical Length (proportion) + Colour Hue (High = dark, Low = light)}
  {Compare proportions across visibility bins}
  { Can accomodate dozons of keys and thousands of values for each key}
  
\insightbox{The relationship between visibility and severity is
\emph{non-monotonic}. Very High visibility ($>10$\,mi) shows the
\textbf{highest} high-severity rate at 31.65\%, while Very Low visibility
($\leq1$\,mi) shows one of the lowest at 16.92\%. Clear, high-visibility roads
(typically highways) produce accidents with longer traffic disruption, whereas
fog accidents on local roads may clear faster. The visibility metric tracks road
type as much as weather condition.}

%% ── H1 Verdict ──────────────────────────────────────────────
\verdict{H1}{Partially Supported}{%
  Rain and Cloudy conditions show the highest severity rates (24.38\% and
  22.43\% respectively after outlier removal). The Rain--Clear gap widened to
  6.53\,pp, representing a stronger signal than pre-cleaning. Snow/Ice (21.45\%)
  has also risen, suggesting the original dataset's extreme duration/distance
  outliers were suppressing adverse-weather severity values. However, Fog/Haze
  still underperforms expectations at 16.93\%, consistent with fog accidents
  being localised and clearing quickly. Weather remains one factor among many.%
}
